\section{Surface Language}

To allow the user to execute surface programs directly we formalise the evaluation of surface terms in
\figref{baseline:surface:eval} and the corresponding notion of surface value. Surface values differ from core
values only in the case of closures, which contain surface expressions other surface values. Later we prove
that directly evaluating a surface program $s$ and desugaring the resulting value is equivalent to desugaring
$s$ to a core program and then evaluating, justifying our presentation of the surface semantics.

\begin{lemma}[Determinism]
   \label{lem:baseline:surface:eval:determinism}
   If $\gamma, s \evalS v$ and $\gamma, s \evalS v'$ then $v = v'$.
\end{lemma}

\subsubsection{Mutual recursion}

The following relation determines the closures associated with a block of mutually recursive functions $g =
\seq{\bind{x}{\mu}}$. (Read $x$ as the function name and $\mu$ as the defining clauses of the function.)

\begin{definition}[Close defs]
   \label{def:baseline:surface:close-defs}
   \figref{baseline:surface:eval} defines the \emph{close defs} relation $\closeDefs$.
\end{definition}

\subsubsection{Piecewise function definitions}

We disallow certain combinations of patterns which are commonly allowed in other languages, such as:%
{\small
\begin{lstlisting}[language=Fluid]
   foo (Cons y ys) (Cons z zs) = ...
   foo  x           Nil        = ...

   bar x ... = ...
   bar y ... = ...
\end{lstlisting}}

\noindent In particular, ``aligned'' variable patterns must have the same name, and variable patterns and any other kind of pattern may only be aligned if we have already been able to distinguish the clauses they are part of. Thus the following is allowed:%
{\small
\begin{lstlisting}[language = Fluid]
   baz (Cons y ys)  x = ...
   baz  Nil         x = ...
\end{lstlisting}}

Let $k$ range over $(\pi, \clause{\pi'}{s})$ pairs. The first component $\pi$ is a stack of subpatterns active
during the processing of a single top-level pattern $p$, initially containing only $p$ and ending up empty.
The second component $\clause{\pi'}{s}$ stores the remaining top-level patterns, and is non-empty only for
curried functions. The following notation is useful for processing piecewise function definitions:

\begin{definition}
   Suppose $k = (\pi, \clause{\pi'}{s})$.
   \begin{enumerate}
      \item Define $p \clauseWith{\cons} k \eqdef (p \cons \pi, \clause{\pi'}{s})$.
      \item Define $\pi^\dagger \clauseWith{\concat} k \eqdef (\pi^\dagger \concat \pi, \clause{\pi'}{s})$.
   \end{enumerate}
   Extend the notation to same-length sequences, so that for example $\seq{p} \clauseWith{\cons} \seq{k} = \seq{p \clauseWith{\cons} k}$.
\end{definition}

\begin{figure}
   {\small
   \begin{minipage}[t]{0.48\textwidth}
   \begin{tabularx}{\textwidth}{rL{2.8cm}L{3cm}}
      &\textbfit{Surface term}&
      \lowlight{$\SExpr\;\Gamma\;A$}
      \\
      $s ::=$
      &
      $\annot{\raw{s}}{\upalpha}$
      \\[2mm]
      &\textbfit{Raw surface term}&
      \\
      $\raw{s} ::=$
      &
      $\exBinaryApp{s}{\oplus}{s'}$
      &
      binary application
      \\
      &
      $\exLetRecPiecewise{g}{s}$
      &
      recursive functions
      \\
      &
      $\exIfThenElse{s}{s_1}{s_2}$
      &
      if
      \\
      &
      $\exMatch{s}{\seq{\clauseUncurried{p}{s}}'}$
      &
      match
      \\
      &
      $\exList{s}{l}$
      &
      non-empty list
      \\
      &
      $\exListEnum{s}{s'}$
      &
      list enum
      \\
      &
      $\exListComp{s}{\seq{q}}$
      &
      list comprehension
      \\
      &
      $\exDo{\seq{d}}$
      &
      do
      \\
      &
      $\exLambda{\upmu}$
      &
      lambda
      \\[2mm]
      &\textbfit{List rest}&
      \lowlight{$\ListRest\;\Gamma\;A$}
      \\
      $\raw{l} ::=$
      &
      $\exListEnd$
      &
      end
      \\
      &
      $\exListNext{s}{l}$
      &
      cons
      \\[2mm]
      &\textbfit{Recursive functions}&
      \\
      $\raw{g} ::=$
      &
      $\seq{\bind{x}{\upmu}}$
      &
      \\
      &\textbfit{Sequential Clauses}&
      \\
      $\raw{\upmu}  ::=$
      &
      $\seq{c}$
      &
      \\
      &\textbfit{Clause}&
      \lowlight{$\Clause\;\Gamma\;A\;B$}
      \\
      $\raw{c} ::=$
      &
      $\clause{\seq{p}}{s}$
      &
      \\
      \end{tabularx}
   \end{minipage}%
   \begin{minipage}[t]{0.5\textwidth}
      \begin{tabularx}{\textwidth}{rL{2.8cm}L{2.9cm}}
      &\textbfit{Pattern}&
      \lowlight{$\Patt\;A\;\Gamma$}
      \\
      $p ::=$
      &
      $\pattVar{x}$
      &
      variable
      \\
      &
      $\pattRecord{\seq{\bind{x}{p}}}$
      &
      record
      \\
      &
      $\pattConstr{c}{\seq{p}}$
      &
      constructor
      \\
      &
      $\pattList{p}{o}$
      &
      non-empty list
      \\[2mm]
      &\textbfit{List rest pattern}&
      \lowlight{$\ListRestPatt\;A\;\Gamma$}
      \\
      $o ::=$
      &
      $\pattListEnd$
      &
      end
      \\
      &
      $\pattListNext{p}{o}$
      &
      cons
      \\[2mm]
      &\textbfit{Qualifier}&
      \\
      $\raw{q} ::=$
      &
      $\qualGuard{s}$
      &
      guard
      \\
      &
      $\qualDeclaration{p}{s}$
      &
      declaration
      \\
      &
      $\qualGenerator{p}{s}$
      &
      generator
      \\[2mm]
      &\textbfit{Do statement}&
      \\
      $\raw{d} ::=$
      &
      $\stmtVoid{s}$
      &
      void
      \\
      &
      $\stmtDeclaration{p}{s}$
      &
      declaration
      \\
      &
      $\stmtBind{p}{s}$
      &
      generator
      \end{tabularx}
   \end{minipage}\\[3mm]
   \begin{minipage}[t]{0.48\textwidth}
      \begin{tabularx}{\textwidth}{rL{2.7cm}L{3cm}}
      &\textbfit{Surface value}&
      \lowlight{$\SVal{A}$}
      \\
      $\raw{u} ::=$
      &
      $\dots$
      &
      $\dots$
      \\
      &
      $\exClosure{\upgamma}{g}{\upmu}$
      &
      closure
      \\[0mm]
      \end{tabularx}
      \end{minipage}
   }
   \caption{Syntax of surface terms}
\end{figure}

\begin{figure}
   {\small \flushleft \tallshadebox{$\upgamma, \seq{s} \evalSugR{\seq{S}} \seq{v}$}%
   \begin{smathpar}
      \inferrule*[
         lab={\ruleName{$\evalSugS$-seq}}
      ]
      {
         \upgamma, s_i \evalSugR{S_i} v_i
         \quad
         (\forall i \numleq \length{\seq{s}})
      }
      {
         \upgamma, \seq{s} \evalSugR{\seq{S}} \seq{v}
      }
      %
      \and
      %
   \end{smathpar}}
   {\small \flushleft \tallshadebox{$\upgamma, s \evalSugR{S} v$}%
   \begin{smathpar}
      \inferrule*[lab={\ruleName{$\evalSugS$-nonempty-list}}]
      {
         \upgamma, s \evalSugR{S} v
         \\
         \upgamma, l \evalSugR{S'} v'
      }
      {
         \upgamma, \exList{s}{l} \evalSugR{\trList{S}{S'}} \cCons(v, v')
      }
      %
      \and
      %
      \inferrule*[
         lab={\ruleName{$\evalSugS$-list-comp}}
      ]{
         \upgamma, \seq{q}, s \evalSugR{\seq{Q},S} v
      }
      {
         \upgamma,  \exListComp{s}{\seq{q}} \evalSugR{\trListComp{S}{Q}} v
      }
      %
      \and
      %
      \inferrule*[
         lab={\ruleName{$\evalSugS$-app}},
         width=3in,
      ]
      {
         \upgamma, (s, s') \evalSugR{(S,S')} (\exClosure{\upgamma_1}{g}{\upmu}, v)
         \\
         \upgamma_1, g \closeDefsS \upgamma_2
         \\
         v, \upmu \matchFwdR{p} \upgamma_3, s^\dag
         \\
         \upgamma_1 \concat \upgamma_2 \concat \upgamma_3, s^\dag \evalSugR{S^\dag} v'
      }
      {
         \upgamma, \exApp{s}{s'} \evalSugR{\trApp{S}{S'}{p}{S^{\dag}}} v'
      }

      %
      \and
      %
      \inferrule*[
         lab={\ruleName{$\evalSugS$-bin-app}}
      ]
      {
         \upgamma, (s_1, s_2) \evalSugR{S_1, S_2} (u_1, u_2)
         \\
         \upgamma, \hat{\oplus}(u_1, u_2) = u
      }
      {
         \upgamma, s_1 \oplus s_1 \evalSugR{\trBinApp{S_1}{S_2}{t}} u
      }
      %
      \and
      %
      \inferrule*[
         lab={\ruleName{$\evalSugS$-if-true}}
      ]
      {
         \upgamma, s \evalSugR{S} \cTrue
         \\
         \upgamma, s_1 \evalSugR{S_1} u
      }
      {
         \upgamma, \exIfThenElse{s}{s_1}{s_2}\evalSugR{\exIfThenElse{S}{S_1}{s_2}} u
      }
      %
      \and
      %
      \inferrule*[
         lab={\ruleName{$\evalSugS$-if-false}}
      ]
      {
         \upgamma, s \evalSugR{S} \cFalse
         \\
         \upgamma, s_2 \evalSugR{S_2} u
      }
      {
         \upgamma, \exIfThenElse{s}{s_{1}}{s_{2}}\evalSugR{\exIfThenElse{S}{s_1}{S2}} u
      }
      %
      \and
      %
      \inferrule*[
         lab={\ruleName{$\evalSugS$-let-rec}}
      ]
      {
         \upgamma, g \closeDefsS \upgamma'
         \\
         \upgamma \concat \upgamma', s \evalSugR{S} v
      }
      {
         \upgamma, \exLetRecPiecewise{g}{s} \evalSugR{\trLetRec{g}{S}} v
      }
      %
      \and
      %
      \inferrule*[
         lab={\ruleName{$\evalSugS$-match-as}}
      ]
      {
         \upgamma, s \evalSugR{S} v
         \\
         v, \seq{\clauseUncurried{p}{s}}' \matchFwdS \upgamma', s_i
         \\
         \upgamma \concat \upgamma', s_i \evalSugR{S'} v'
      }
      {
         \upgamma, \exMatch{s}{\seq{\clauseUncurried{p}{s}}'} \evalSugR{\trMatchAs{S}{\clauseUncurried{p'}{S'}}} v'
      }
      %
      \and
      %
      \inferrule*[
         lab={\ruleName{$\evalSugS$-enum}}
      ]
      {
         \upgamma, \exApp{\exApp{\varEnumFromTo}{s_1}}{s_2} \evalSugR{\trSurfaceApp{(\trSurfaceApp{\varEnumFromTo}{S_1}{\exClosure{\varepsilon}{\varepsilon}{\clause{y}{S'}}})}{S_2}{S^{\dag}}} v
      }
      {
         \upgamma, \exListEnum{s_1}{s_2} \evalSugR{\trEnum{S_1}{S_2}} v
      }
   \end{smathpar}}
   {\small \flushleft \tallshadebox{$\seq{q}, s \evalSugR{\seq{Q},S} v$}
   \begin{smathpar}
      \inferrule*[lab={\ruleName{$\evalSugS$-list-comp-done}}]
      {
         \upgamma, s \evalSugR{S} u
      }
      {
         \upgamma, \exListComp{s}{\seqEmpty} \evalSugR{\seqEmpty, S} \cCons(u,\cNil)
      }
      %
      \and
      %
      \inferrule*[
         lab={\ruleName{$\evalSugS$-list-comp-decl}}
      ]
      {
         \upgamma, s' \evalSugR{S'} u'
         \\
         p, u' \matchS \upgamma'
         \\
         (\upgamma \cons \upgamma'), \seq{q},s \evalSugR{\seq{Q},S} u
      }
      {
         \upgamma, \qualDeclaration{p}{s'} \cons \seq{q}, s \evalSugR{\trQualDecl{p}{S'} \cons \seq{Q}, S} u
      }
      %
      \and
      %
      \inferrule*[
         lab={\ruleName{$\evalSugS$-list-comp-gen}}
      ]
      {
         \upgamma, \exApp{\exApp{\varConcatMap}{\uplambda \{p = \exListComp{s'}{\seq{q}} \}}}{s} \evalSugS u
      }
      {
         \upgamma, \qualGenerator{p}{s'} \cons \seq{q}, s \evalSugR{\trQualGen{p}{S'}\cons\seq{Q}, S} u
      }
      %
      \and
      %
      \inferrule*[
         lab={\ruleName{$\evalSugS$-list-comp-guard-false}}
      ]
      {
         \upgamma, s \evalSugR{S} \cFalse
      }
      {
         \upgamma, \qualGuard{s} \cons \seq{q}, s' \evalSugR{\trQualIfFalse{S}{q}} []
      }
      %
      \and
      %
      \inferrule*[
         lab={\ruleName{$\evalSugS$-list-comp-guard-true}},
      ]
      {
         \upgamma, \seq{q}, s' \evalSugR{\seq{Q},S} v
         \\
         \upgamma, s \evalSugR{S} \cTrue
      }
      {
         \upgamma, \qualGuard{s} \cons \seq{q}, s' \evalSugR{\trQualIfTrue{S} \cons \seq{Q}, S} v
      }
   \end{smathpar}}
   \caption{Surface semantics: evaluation}
\end{figure}
\begin{figure}
   \ContinuedFloat
   {\small \flushleft \tallshadebox{$\upgamma, r \evalSugS v$}
   \begin{smathpar}
      \inferrule*[
         lab={\ruleName{$\evalSugS$-list-rest-end}}
      ]
      {
         \strut
      }
      {
         \upgamma, \exListEnd \evalSugR{\trListEnd} \cNil
      }
      %
      \and
      %
      \inferrule*[
         lab={\ruleName{$\evalSugS$-list-rest-next}}
      ]
      {
         \upgamma, s \evalSugR{S} v
         \\
         \upgamma, l \evalSugR{S'} v'
      }
      {
         \upgamma, \exListNext{s}{l} \evalSugR{\trListNext{S}{S'}} \cCons(v, v')
      }
   \end{smathpar}
   }
   \caption{Surface semantics: evaluation}
\end{figure}
\begin{figure}
   {\small, \flushleft \shadebox{$\upgamma, g \closeDefsS \upgamma'$}
      \begin{smathpar}
         \inferrule*
         {
            g = \seq{\bind{x}{\mu}}
            \\
            v_i = \exClosure{\upgamma}{g}{\upmu_i}
            \quad
            \forall i \le \length{\seq{x}}
         }
         {
            \upgamma, g \closeDefsS \set{\seq{\bind{x}{v}}}
         }
      \end{smathpar}}
   \caption{Surface semantics: recursive definitions}
\end{figure}

\begin{figure}
   {\small \flushleft \shadebox{$\seq{v}, \seq{(\pi, k)} \match \gamma, s$}
   \begin{smathpar}
      \inferrule*[
        lab={\ruleName{$\match$-done-expr}}
      ]
      {
         \strut
      }
      {
         \seqEmpty, (\seqEmpty, \clause{\seqEmpty}{s}) \match \seqEmpty, s
      }
      %
      \and
      %
      \inferrule*[
        lab={\ruleName{$\match$-done-lambda}}
      ]
      {
         \strut
      }
      {
         \seqEmpty, (\seqRange{(\seqEmpty, \clause{p_1 \cons \pi_1}{s_1})}
                              {(\seqEmpty, \clause{p_j \cons \pi_j}{s_j})})
         \match
         \seqEmpty, \exFun{(\seq{\clause{p \cons \pi}{s}})}
      }
   \end{smathpar}}
   \\[3mm]
   {\small \flushleft \shadebox{$v, \mu \match \gamma, s$}
   \begin{smathpar}
      \inferrule*[
         lab={\ruleName{$\match$-uncurried}}
      ]
      {
         v, p_{i} \match \gamma
         \\
         (\forall i' \numleq j.i' \ne i \implies v, p_{i'}\;\matchNeg)
      }
      {
         v, (\seqRange{\clauseUncurried{p_1}{s_1}}{\clauseUncurried{p_j}{s_j}}) \match \gamma, s_i
      }
      %
      \and
      %
      \inferrule*[
         lab={\ruleName{$\match$-curried}},
         right={$\mu \neq \seqEmpty$}
      ]
      {
         v, p \match \gamma
         \\
         \mu \eqdef (\clause{\seq{p_i}}{s_i} \mid \forall i \numleq j.p_i = p)
         \\
         (\forall i \numleq j.p_i \neq p \implies v, p_i\;\matchNeg)
      }
      {
         v, (\seqRange{\clause{(p_1 \cons \seq{p_1})}{s_1}}{\clause{(p_j \cons \seq{p_j})}{s_j}})
         \match
         \gamma,
         \exFun{\mu}
      }
   \end{smathpar}}
   \caption{Surface semantics: pattern matching}
\end{figure}

\begin{figure}
   {\small \flushleft \shadebox{$k, s^\dagger \orElse \seq{k}$}
   \begin{smathpar}
      \inferrule*[
         lab={\ruleName{$\orElse$-done}}
      ]
      {
      }
      {
      }
   \end{smathpar}}
   \caption{Completing a clause with a default value $s'$}
\end{figure}

